\section{Introduzione}\label{sec:intro}

Ospitare le Olimpiadi è comunemente associato a un vantaggio competitivo, per
svariati motivi: i fattori geografici includono lunghi viaggi e (estrema)
diversità in quanto a cultura, ora (jetlag), clima; quelli sportivi comprendono
pubblico più caloroso, serenità di giocare in casa, possibili influenze sugli
arbitri, maggiori investimenti per non sfigurare, possibilità di scegliere degli
eventi da aggiungere.\\
Giornali e discussioni informali sono supportati anche da veri studi, nel
mostrare un incremento di medaglie quando si ospiti (\textit{home advantage}).
Tuttavia, osservazioni del genere possono solo mostrare correlazione tra i due
fenomeni, piuttosto che causalità.\\
Dunque, il nostro scopo qui, dopo aver constatato una correlazione, è
verificare, nelle Olimpiadi estive, se un eventuale boost prestazionale sia
influenzato principalmente dai fattori di ospitalità, invece che da altri.